% ----------------------------------------------------------------------
% common text abbreviations
\newcommand{\eg}{\mbox{\it e.g.}}
\newcommand{\ie}{\mbox{\it i.e.}}
\newcommand{\etal}{\mbox{\it et al.}}
\newcommand{\apriori}{{\it a priori}}
\newcommand{\aposteriori}{{\it a posteriori}}
\newcommand{\wloss}{without loss of generality}
\newcommand{\Wloss}{Without loss of generality}

% ----------------------------------------------------------------------
% common symbols

% implication
\newcommand{\imp}{\Rightarrow}                  
\newcommand{\iimp}{\Longrightarrow}       
\renewcommand{\implies}{$\imp\ $}		
\newcommand{\pmi}{\Leftarrow}		

% spacing
\newcommand{\sep}{\hspace{2em}} 
\newcommand{\ssep}{\hspace{1em}}
\newcommand{\void}[1]{}				% comment out the argument
%\newcommand{\todo}[1]{[#1]\marginpar{\mbox{\Huge !}}}

%Here are some user-defined notations

\newcommand{\RR}{\mathbf R}
\newcommand{\CC}{\mathbf C}
\newcommand{\ZZ}{\mathbf Z}
\newcommand{\ZZn}[1]{\ZZ/{#1}\ZZ}
\newcommand{\QQ}{\mathbf Q}
\newcommand{\rr}{\mathbb R}
\newcommand{\cc}{\mathbb C}
\newcommand{\zz}{\mathbb Z}
\newcommand{\ff}{\mathbb{F}}
\newcommand{\kk}{\mathbf k}
\newcommand{\zzn}[1]{\zz/{#1}\zz}
\newcommand{\qq}{\mathbb Q}
\newcommand{\nn}{\mathbb N}
\newcommand{\calM}{\mathcal{M}}
\newcommand{\calL}{\mathcal{L}}
\newcommand{\dd}{{\rm d}}
\newcommand{\set}[1]{\mathcal{#1}}
\newcommand{\scr}[1]{\mathscr{#1}}


\renewcommand{\leq}{\leqslant}		
\renewcommand{\geq}{\geqslant}	
\renewcommand{\preceq}{\preccurlyeq}
\renewcommand{\succeq}{\succcurlyeq}

% fonts
\newcommand{\bm}{\boldsymbol}
\newcommand{\wt}{\widetilde}

\newcommand{\bma}{{\bm a}}
\newcommand{\bmb}{{\bm b}}
\newcommand{\bmc}{{\bm c}}
\newcommand{\bmd}{{\bm d}}
\newcommand{\bme}{{\bm e}}
\newcommand{\bmx}{{\bm x}}
\newcommand{\bmy}{{\bm y}}
\newcommand{\bmz}{{\bm z}}
\newcommand{\bmu}{{\bm u}}
\newcommand{\bmv}{{\bm v}}
\newcommand{\bmw}{{\bm w}}
\newcommand{\bms}{{\bm s}}
\newcommand{\bmt}{{\bm t}}

\newcommand{\bmmu}{{\bm\mu}}
\newcommand{\bmlambda}{{\bm\lambda}}

\newcommand{\bb}{{\bm b}}
\newcommand{\bbb}{{\bm b}}
\newcommand{\bbx}{{\bm x}}
\newcommand{\bcc}{{\bm c}}
\newcommand{\xx}{{\bm x}}
\newcommand{\yy}{{\bm y}}
\newcommand{\uu}{{\bm u}}
\newcommand{\vv}{{\bm v}}
\newcommand{\ww}{{\bm w}}


\newcommand{\calP}{\mathcal P}
\newcommand{\calS}{\mathcal S}
\newcommand{\calH}{\mathcal H}
\newcommand{\cat}{\mathcal}




% operatornames

\newcommand{\inner}[1]{\left\langle#1\right\rangle}
\newcommand{\aff}{\operatorname{aff}}
\newcommand{\ri}{\operatorname{ri}}
\newcommand{\cl}{\operatorname{cl}}
\newcommand{\bd}{\operatorname{bd}}
\newcommand{\dist}{\operatorname{dist}}
\newcommand{\Spec}{\operatorname{Spec}}
\newcommand{\rec}{\operatorname{rec}}
\newcommand{\cone}{\operatorname{cone}}
\newcommand{\conv}{\operatorname{conv}}
\newcommand{\im}{\operatorname{im}}
\newcommand{\diam}{\operatorname{diam}}

\renewcommand{\subset}{\subseteq}
\newcommand{\superset}{\supseteq}
\newcommand{\supersetneq}{\supsetneq}

% shortcuts 
\newcommand{\bmat}[1]{\begin{bmatrix} #1 \end{bmatrix}}
\newcommand{\rast}{^\ast}
\newcommand{\e}{\operatorname{e}}
\newcommand{\iu}{\operatorname{i}}
\newcommand{\pp}{\mathbf{P}}
\newcommand{\ppn}{\pp^{n}}
\newcommand{\prob}{\mathbb{P}}
\newcommand{\ee}{\mathbb{E}}
\newcommand{\rise}[1]{^{(#1)}}
\newcommand{\mr}{\operatorname{mr}}
\newcommand{\paints}[2]{#1\to #2}
\newcommand{\rank}{\operatorname{rank}}
\newcommand{\corank}{\operatorname{corank}}
\newcommand{\supp}{\operatorname{supp}}
\newcommand{\vol}{\textsc{Vol}}
\newcommand{\Span}{\operatorname{span}}

\newcommand{\spectralnorm}[1]{{\left\vert\kern-0.25ex\left\vert\kern-0.25ex\left\vert #1 
    \right\vert\kern-0.25ex\right\vert\kern-0.25ex\right\vert}}
    

%\setcounter{chapter}{-1}

\numberwithin{equation}{chapter}
% ----------------------------------------------------------------------
% theorem definitions
% numbered by default: 
% 	theorem, lemma, proposition, corollary, example, principle
% unnumbered by default: 
%	definition, openproblem, remark, remarks, convention
%
% get unnumbered versions by un-lemma etc., numbered by n-definition etc. 



\theoremstyle{plain}
\newtheorem{n-lemma}[equation]{Lemma}
\newtheorem{n-theorem}[equation]{Theorem}
\newtheorem{n-proposition}[equation]{Proposition}
\newtheorem{n-corollary}[equation]{Corollary} 
\newtheorem{n-principle}[equation]{Principle}
\newtheorem{n-claim}[equation]{Claim}
\newtheorem{n-assumption}[equation]{Assumption}

\newtheorem*{un-lemma}{Lemma}
\newtheorem*{un-theorem}{Theorem}
\newtheorem*{un-proposition}{Proposition}
\newtheorem*{un-corollary}{Corollary} 
\newtheorem*{un-principle}{Principle}
\newtheorem*{un-claim}{Claim}

\theoremstyle{definition}
\newtheorem{n-definition}[equation]{Definition}
\newtheorem{n-condition}[equation]{Condition}
\newtheorem{n-openproblem}[equation]{Open Problem}
\newtheorem{bn-remark}[equation]{Remark}		% bold remark

\newtheorem*{un-definition}{Definition}
\newtheorem*{un-openproblem}{Open Problem}
\newtheorem*{bun-remark}{Remark}		% bold unnumbered remark
\newtheorem*{un-problem}{Problem}
\newtheorem{n-problem}[equation]{Problem}

\theoremstyle{remark}
\newtheorem{n-example}[equation]{Example}
\newtheorem{n-remark}[equation]{Remark}
\newtheorem{n-remarks}[equation]{Remarks}
\newtheorem{n-convention}[equation]{Convention}
\newtheorem*{un-example}{Example}
\newtheorem*{un-remark}{Remark}
\newtheorem*{un-remarks}{Remarks}
\newtheorem*{un-convention}{Convention}

% defaults numbered/unnumbered
\newenvironment{theorem}{\begin{mdframed}[linewidth=1pt,            
        innerleftmargin=10pt,
        innertopmargin=8pt,
        skipabove=10pt] 
\vspace{-\topsep}\begin{n-theorem}}{\end{n-theorem}\end{mdframed}}
\newenvironment{lemma}{\begin{mdframed}[linewidth=1pt,            
        innerleftmargin=10pt,
        innertopmargin=8pt,
        skipabove=10pt] 
\vspace{-\topsep}\begin{n-lemma}}{\end{n-lemma}\end{mdframed}}
\newenvironment{proposition}{\begin{mdframed}[linewidth=1pt,           
        innerleftmargin=10pt,
        innertopmargin=8pt,
        skipabove=10pt] 
\vspace{-\topsep}\begin{n-proposition}}{\end{n-proposition}\end{mdframed}}
\newenvironment{corollary}{\begin{mdframed}[linewidth=1pt,            
        innerleftmargin=10pt,
        innertopmargin=8pt,
        skipabove=10pt] 
\vspace{-\topsep}\begin{n-corollary}}{\end{n-corollary}\end{mdframed}}
\newenvironment{definition}{\begin{mdframed}[linewidth=1pt,            
        innerleftmargin=10pt,
        innertopmargin=8pt,
        skipabove=10pt] 
\vspace{-\topsep}\begin{n-definition}}{\end{n-definition}\end{mdframed}}
\newenvironment{condition}{\begin{n-condition}}{\end{n-condition}}
\newenvironment{assumption}{\begin{n-assumption}}{\end{n-assumption}}
\newenvironment{example}{\begin{mdframed}[linewidth=1pt,            
        innerleftmargin=10pt,
        innertopmargin=8pt,
        skipabove=10pt] 
\begin{n-example}}{\end{n-example}\end{mdframed}}
\newenvironment{openproblem}{\begin{un-openproblem}}{\end{un-openproblem}}
\newenvironment{principle}{\begin{n-principle}}{\end{n-principle}}
\newenvironment{remark}{\begin{mdframed}[linewidth=1pt,            
        innerleftmargin=10pt,
        innertopmargin=8pt,
        skipabove=10pt] 
\vspace{-\topsep}\begin{n-remark}}{\end{n-remark}\end{mdframed}}
\newenvironment{remarks}{\begin{un-remarks}}{\end{un-remarks}}
\newenvironment{convention}{\begin{un-convention}}{\end{un-convention}}
\newenvironment{claim}{\begin{un-claim}}{\end{un-claim}}
\newenvironment{problem}{\begin{n-problem}}{\end{n-problem}}


%\newaliascnt{equation}{n-lemma}
%\aliascntresetthe{equation}

% proofs
\renewenvironment{proof}
{\noindent{\em Proof:}\hspace{.5em}\textrm}
{\hfill\qedsymbol

\vspace{1em}}




\newenvironment{proofof}[1]			% loose proofs
	{\begin{trivlist}\item[]{\em  Proof of #1:\hspace{.5em}}\rm}
% 	{\end{trivlist}}
	{\hfill\qedsymbol
	\end{trivlist}}
\newenvironment{proofsketch}			% proof sketch
	{\begin{trivlist}\item[]{\em  Proof sketch:\hspace{.5em}}\rm}
	{\hfill\qedsymbol
	\end{trivlist}}

\renewcommand{\qed}{} 		% turn off automatic AMS-style \qed'ing.
\newcommand{\eotsymbol}{{\large $\Box$}} 	% can be redefined
\newcommand{\eot}{{\hspace*{\fill}\eotsymbol}}
\newcommand{\eottwo}{\vspace{-5ex}\par\eot}	% \eot after math display

% ----------------------------------------------------------------------
% case distinctions: label is flush left. 
% caselist causes no indentation. May use \case for automatic numbering
% itemlist causes indentation. 
% nienumerate is like enumerate but flush left

\newcounter{casecounter}
\newlength{\normalparindent}
\setlength{\normalparindent}{\parindent}
\newenvironment{caselist}
	{\begin{list}{{\bf Case \arabic{casecounter}:}}
	 {\usecounter{casecounter}
		\setlength{\leftmargin}{0ex}
		\setlength{\labelwidth}{0ex} 		 
		\setlength{\labelsep}{2ex}
		\setlength{\itemindent}{\labelsep}	
		\setlength{\parsep}{0ex}
		\setlength{\listparindent}{\normalparindent}
	}}	
	{\end{list}}
\newcommand{\case}{\item}
\newcommand{\caseimp}{\case[$\imp$:]}
\newcommand{\casepmi}{\case[$\pmi$:]}
\newcommand{\impl}[2]{$\text{#1}\imp\text{#2}$}
\newcommand{\caseimpl}[2]{\case[\impl{#1}{#2}:]}
\newenvironment{itemlist}{\begin{list}{\undefinedcase}
 	{	\setlength{\leftmargin}{2.5ex}
		\setlength{\labelwidth}{1.5ex}
 		\setlength{\labelsep}{1ex}
		\setlength{\itemindent}{0ex}}}
 	{\end{list}}

\newcounter{nicounter}
\newenvironment{nienumerate} % enumerate but no indentation
	{\begin{list}{\arabic{nicounter}.}
{\usecounter{nicounter} 		 \setlength{\leftmargin}{0ex}
\setlength{\labelwidth}{0ex} 		 \setlength{\labelsep}{1ex}
\setlength{\itemindent}{\labelsep} 	}	} 	{\end{list}}



